\section{估计与回测方法细节}
\label{app:methods}
\subsection{绩效指标与日期口径}
\label{app:metrics}
\subsubsection{自然周与持有期}
自然周收益由该周所有日收益复利得到。调仓后持有期则从决策日开始，到下一次决策前结束。由于主模型在周末调仓，二者的日期集合通常不同。将自然周收益直接附在本周末目标权重旁边，并将其称为该目标的后续收益，会产生解释错误。

\subsubsection{年化收益与波动}
日净值由 $V_t=\prod_{s\leq t}(1+r_s^{\rm net})$ 构造，年化收益按净值首末比值计算：
\[
R_{\rm ann}=\left(\frac{V_N}{V_1}\right)^{\annualizationDays/N}-1.
\]
首个评价日净值 $V_1$ 为基期净值，波动和夏普比率使用相邻净值变化计算。

设由净值恢复的日收益为 $u_t$，年化波动为
\[
\sigma_{\rm ann}=\operatorname{sd}(u_t)\sqrt{\annualizationDays}.
\]
标准差采用样本自由度。年化系数是研究约定，不意味着每个自然年实际交易日数均恰好相同。部分年度的年化收益也不同于截至该日的实际累计收益，年度表格会明确标注。

\subsubsection{夏普、Sortino 与回撤}
夏普比率计算如下\footnote{本文统一设定无风险利率 $r_f=0$，夏普比率以日净收益相对于零收益基准计算，Sortino 比率的最低可接受收益亦设为零。该设定为五模型比较的统计口径，不表示实际无风险资产收益为零。}：
\[
S=\frac{\operatorname{mean}(u_t)}{\operatorname{sd}(u_t)}\sqrt{\annualizationDays}.
\]
该分子为算术平均收益，不是几何年化收益。因此，不能用表中的净年化收益除以年化波动并要求得到完全相同的夏普。Sortino 用负收益的二阶矩替代标准差，正收益对应下行项为零，均值仍基于完整收益序列。

最大回撤定义为
\[
D_t=\frac{V_t}{\max_{s\leq t}V_s}-1,\qquad
D_{\max}=\min_t D_t.
\]
表格以负数表示最深下跌。回撤关注路径及先后顺序，与只观察收益分布的波动不同。同样的日收益集合若排列次序不同，最终累计收益可以相同，而最大回撤可能不同。

\subsubsection{尾部损失与换手汇总}
令日损失 $L_t=-r_t^{\rm net}$。项目先计算损失的 95\% 样本分位数，再对不低于该分位数的损失取平均，作为报告的历史 CVaR。有限样本中的分位数插值及边界并列值，会使入选观察数量不必恰好等于样本的 5\%。正文的 CVaR 均沿用这一实现。

月均换手为全部调仓换手之和除以评价区间覆盖的自然月数，而非每次交易的平均值。年化收益成本拖累定义为同一路径毛、净年化收益之差，包含复利影响，不等于简单将费率乘以自然年换手。分别定义这些量，可以防止在图表解释中多扣一次费用或混用年化方式。


\subsection{对照模型的分配规则}
\label{app:benchmarks}
\subsubsection{Equal Weight}
Equal Weight 在每个周度决策日，将资本平均分配给当期合格资产。若合格资产为 $n_t$ 个，则目标为 $w_i=1/n_t$。这种方法不需要收益预测，也不需要将协方差矩阵输入分配公式。资产的上市及准入仍由历史窗口决定，因此等权对象随历史可用性变化。

等权的优点在于规则透明，结果差异容易与资产池结构联系。在本研究中，多只权益和商品工具并存，等资本会使高波动资产承担较多组合风险。等权收益较高时，不能据此推断其风险控制更好；同样，较大回撤也不意味着等权失去参照价值。它直接呈现持有合格资产集合本身所带来的资本暴露。

\subsubsection{60/40 Benchmark}
60/40 Benchmark 按代码的资产分类，将 60\% 资本分给权益组，40\% 分给债券组，并在各组内等权。设合格权益与债券数量分别为 $n_E,n_B$，则
\[
w_i=\begin{cases}
0.60/n_E,&i\in E,\\
0.40/n_B,&i\in B,\\
0,&\text{其他资产}.
\end{cases}
\]
该对照是固定资本预算，不是风险预算。两组的实际风险贡献会随着波动与相关性变化，权益风险可能明显超过其资本比例。若某一期缺少必要资产组，则程序依既有规则处理，而不使用未来资产补足组合。
